\section{Designs That Did Not Identify Their Target}\label{app:designs}

This appendix reports three designs that were run and did not identify what they were meant to, so that the reader
can see every test behind the argument.

\paragraph{Gender-specific labor-demand shocks at a birth.} To separate the mother's re-employment friction from the
husband's risk, I built shift-share (Bartik) shocks to employment in women's and men's industries by state from the CPS
Outgoing Rotation Groups (1979--2019, 16 sectors). In a split-sample test, the women's shock does not predict women's
state employment growth ($-0.25$, s.e.\ 0.18), while the men's shock predicts men's (0.24, s.e.\ 0.11). With only 16
sectors, the variation left after state and year fixed effects is small, and the leave-one-out construction adds
mechanical noise from small CPS cells. The women's shock is therefore not used. Code: \texttt{IGE/8a.StateShocks.py},
\texttt{IGE/8.FirstStage.py}.

\paragraph{The Family and Medical Leave Act.} I compared births to NLSY79 mothers eligible for job-protected leave
(employer of 50 or more at the location, 12 months of tenure, 1,250 hours) before and after August 1993. The
difference-in-differences estimate on returning to the pre-birth employer within 12 weeks is $-0.073$ (s.e.\ 0.034), and
the estimates at 26 and 52 weeks are not significant. The design is weak: eligible mothers were already near the
ceiling of returning (0.88--0.93), the 1994 redesign of the NLSY79 coincides with the post period and changes the
measured frequency of breaks, and returning within 12 weeks mechanically misses mothers who took the full 12 weeks of
leave. Code: \texttt{IGE/12.FMLA.py}.

\paragraph{State-by-sector shocks to husbands.} Employment growth in the husband's sector within his state has no
first stage on the husband's hours or earnings (F below 1), even with split-sample instruments, so only national sector
shocks are used in Appendix~\ref{sec:birthshocks}. Code: \texttt{IGE/18.HusbandIndustryShocks.py}.
